% Options for packages loaded elsewhere
\PassOptionsToPackage{unicode}{hyperref}
\PassOptionsToPackage{hyphens}{url}
\PassOptionsToPackage{dvipsnames,svgnames,x11names}{xcolor}
\documentclass[
  10pt,
]{article}
\usepackage{xcolor}
\usepackage[margin=2.4cm]{geometry}
\usepackage{amsmath,amssymb}
\setcounter{secnumdepth}{-\maxdimen} % remove section numbering
\usepackage{iftex}
\ifPDFTeX
  \usepackage[T1]{fontenc}
  \usepackage[utf8]{inputenc}
  \usepackage{textcomp} % provide euro and other symbols
\else % if luatex or xetex
  \usepackage{unicode-math} % this also loads fontspec
  \defaultfontfeatures{Scale=MatchLowercase}
  \defaultfontfeatures[\rmfamily]{Ligatures=TeX,Scale=1}
\fi
\usepackage{lmodern}
\ifPDFTeX\else
  % xetex/luatex font selection
  \setmainfont[]{Libertinus Serif}
  \setmonofont[Scale=0.85]{Latin Modern Mono}
  \setmathfont[]{Latin Modern Math}
\fi
% Use upquote if available, for straight quotes in verbatim environments
\IfFileExists{upquote.sty}{\usepackage{upquote}}{}
\IfFileExists{microtype.sty}{% use microtype if available
  \usepackage[]{microtype}
  \UseMicrotypeSet[protrusion]{basicmath} % disable protrusion for tt fonts
}{}
\makeatletter
\@ifundefined{KOMAClassName}{% if non-KOMA class
  \IfFileExists{parskip.sty}{%
    \usepackage{parskip}
  }{% else
    \setlength{\parindent}{0pt}
    \setlength{\parskip}{6pt plus 2pt minus 1pt}}
}{% if KOMA class
  \KOMAoptions{parskip=half}}
\makeatother
\usepackage{longtable,booktabs,array}
\newcounter{none} % for unnumbered tables
\usepackage{calc} % for calculating minipage widths
% Correct order of tables after \paragraph or \subparagraph
\usepackage{etoolbox}
\makeatletter
\patchcmd\longtable{\par}{\if@noskipsec\mbox{}\fi\par}{}{}
\makeatother
% Allow footnotes in longtable head/foot
\IfFileExists{footnotehyper.sty}{\usepackage{footnotehyper}}{\usepackage{footnote}}
\makesavenoteenv{longtable}
\setlength{\emergencystretch}{3em} % prevent overfull lines
\providecommand{\tightlist}{%
  \setlength{\itemsep}{0pt}\setlength{\parskip}{0pt}}

\providecommand{\E}{\mathbb{E}}
\providecommand{\N}{\mathbb{N}}
\providecommand{\Z}{\mathbb{Z}}
\providecommand{\R}{\mathbb{R}}
\providecommand{\eps}{\varepsilon}
\providecommand{\ceil}[1]{\left\lceil #1\right\rceil}
\providecommand{\floor}[1]{\left\lfloor #1\right\rfloor}
\AtBeginDocument{\let\setminus\smallsetminus}
\usepackage[most]{tcolorbox}
\newtcolorbox{warning}{colback=red!5,colframe=red!65!black,boxrule=0.8pt,arc=2pt,left=6pt,right=6pt,top=5pt,bottom=5pt}
\usepackage{etoolbox}
\AtBeginEnvironment{longtable}{\small}
\setlength{\emergencystretch}{3em}
\usepackage{fancyhdr}
\pagestyle{fancy}
\fancyhf{}
\fancyhead[L]{\small\bfseries UNREFEREED GPT-6 ASTRA OUTPUT}
\fancyhead[R]{\small\thepage}
\renewcommand{\headrulewidth}{0.4pt}
\usepackage{bookmark}
\IfFileExists{xurl.sty}{\usepackage{xurl}}{} % add URL line breaks if available
\urlstyle{same}
\hypersetup{
  pdftitle={Informal Conjecture on Property (*) for d = o(n/log n)},
  pdfauthor={Writeup: gpt-6-astra (single model pass)},
  colorlinks=true,
  linkcolor={blue!50!black},
  filecolor={Maroon},
  citecolor={Blue},
  urlcolor={blue!50!black},
  pdfcreator={LaTeX via pandoc}}

\title{Informal Conjecture on Property (*) for d = o(n/log n)}
\usepackage{etoolbox}
\makeatletter
\providecommand{\subtitle}[1]{% add subtitle to \maketitle
  \apptocmd{\@title}{\par {\large #1 \par}}{}{}
}
\makeatother
\subtitle{Unrefereed candidate proof by GPT-6 Astra}
\author{Writeup: \texttt{gpt-6-astra} (single model pass)}
\date{Catalog id \texttt{2207.13651\_\_00} --- generated 2026-09-17}

\begin{document}
\maketitle

\begin{warning}

\textbf{UNREFEREED MODEL OUTPUT.} This document was selected solely
because GPT-6 Astra labelled its own result
\texttt{would\_publish:\ true} and returned \texttt{proved} or
\texttt{disproved}. It has not passed the adversarial LLM referee used
for the earlier Sol campaign, has not been checked by a human
mathematician, and has not been checked for novelty. Treat every
mathematical and bibliographic claim below as unverified.

\end{warning}

\section{Summary and provenance}\label{summary-and-provenance}

{\def\LTcaptype{none} % do not increment counter
\begin{longtable}[]{@{}
  >{\raggedright\arraybackslash}p{(\linewidth - 2\tabcolsep) * \real{0.5000}}
  >{\raggedright\arraybackslash}p{(\linewidth - 2\tabcolsep) * \real{0.5000}}@{}}
\toprule\noalign{}
\endhead
\bottomrule\noalign{}
\endlastfoot
Catalog id & \texttt{2207.13651\_\_00} \\
Catalog entry &
\href{https://graph-theory-ai.github.io/graph-conjectures/arxiv/2207.13651__00/}{Informal
Conjecture on Property (*) for d = o(n/log n)} \\
Source corpus & arxiv \\
Campaign leg & selected second attempts (\texttt{attacks\_retry}) \\
Source paper / entry & On random irregular subgraphs \\
Model verdict & \textbf{proved} (confidence: high) \\
Model's one-line claim & A half-interval decomposition of the vertex
labels, followed by read-d Hölder bounds for uniform spacings, proves
Property (*) throughout d = o(n/log n). \\
Model & \texttt{gpt-6-astra}, reasoning effort \texttt{max},
\texttt{mode=pro}, flex service tier \\
Original artifact &
\texttt{attacks\_retry/2207.13651\_\_00/output.md} \\
Independent review & \textbf{None} \\
Model's caveats & Literature priority has not been independently
checked; graphs are finite, simple, and undirected. \\
\end{longtable}
}

\section{Problem statement}\label{problem-statement}

\emph{No extracted statement was stored in the artifact metadata.}

\section{Astra writeup}\label{astra-writeup}

\emph{The text below is the model output verbatim, apart from moving its
machine-readable verdict block into the summary above and shifting
Markdown heading levels for this document. Mechanical TeX defects and
equation tags were normalized where needed for compilation.}

\subsection{\texorpdfstring{Property \((*)\) in the conjectured
range}{Property (*) in the conjectured range}}\label{property-in-the-conjectured-range}

\subsubsection{1. Statement and the key
observation}\label{statement-and-the-key-observation}

Let \(G\) be a \(d\)-regular graph on \(n\) vertices. Independently
assign \(X_v\sim\mathrm{Unif}[0,1]\), and retain \(uv\in E(G)\) exactly
when \(X_u+X_v\ge1\). Write \[
m_k=\bigl|\{v:\deg_H(v)=k\}\bigr|,
\qquad
\mu=\frac{n}{d+1}.
\]

We prove the following.

\paragraph{Theorem}\label{theorem}

For every sequence of \(d\)-regular graphs with \[
d=o(n/\log n),
\] with probability tending to \(1\), \[
\max_{0\le k\le d}|m_k-\mu|=o(\mu).
\] The probability estimates are uniform over the choice of \(G\).

More quantitatively, there is an absolute constant \(C\) such that, for
all sufficiently large \(n\), whenever \(d\ge64\log n\), \[
\Pr\left(
\max_{0\le k\le d}|m_k-\mu|
>
C\left(\sqrt{\frac nd\log n}+\log n\right)
\right)
\le Cn^{-3}.
\qquad\text{(1)}
\]

The useful idea is to record whether each label belongs to the lower or
upper half of \([0,1]\). Conditional on these bits:

\begin{itemize}
\tightlist
\item
  every lower--lower edge is absent;
\item
  every upper--upper edge is present;
\item
  only edges across the resulting bipartition remain random.
\end{itemize}

Thus one can use the independent-side argument from the supplied earlier
attempt \textbf{without properly coloring \(G\)}. The additional random
conditional mean is a sum of bounded local functions of independent
bits, and a read-\((d+1)\) Hölder inequality controls it.

We verify all these steps below.

\begin{center}\rule{0.5\linewidth}{0.5pt}\end{center}

\subsubsection{2. Two elementary tools}\label{two-elementary-tools}

\paragraph{2.1. Product-space Hölder
inequality}\label{product-space-huxf6lder-inequality}

Suppose \(Z_1,\ldots,Z_N\) are independent, and each nonnegative
function \(f_i\) depends on some subset of these coordinates. If every
coordinate occurs in at most \(R\) functions, then \[
\mathbb E\prod_i f_i
\le
\prod_i\left(\mathbb E f_i^R\right)^{1/R}.
\qquad\text{(2)}
\]

For completeness, integrate one coordinate at a time. At a coordinate
occurring in \(q\le R\) factors, apply Hölder with exponent \(R\) to
those \(q\) factors, adding a constant factor when \(q<R\). The
resulting factors have the form \[
\left(\mathbb E_{Z_j} f_i^R\right)^{1/R}.
\] Repeating over the remaining coordinates gives (2).

In particular, if \(Y_i\) are functions with this read-\(R\) property,
then \[
\log\mathbb E e^{\lambda\sum_i(Y_i-\mathbb EY_i)}
\le
\frac1R\sum_i
\log\mathbb E e^{\lambda R(Y_i-\mathbb EY_i)}.
\qquad\text{(3)}
\]

\paragraph{2.2. Uniform spacings}\label{uniform-spacings}

Let \(B_r\) be any specified spacing between \(r\) independent uniform
points in \([0,1]\), including the two endpoint spacings. Then \[
B_r\sim\operatorname{Beta}(1,r),
\qquad
\mathbb E B_r=\frac1{r+1},
\] and, for every integer \(j\ge1\), \[
\mathbb E B_r^j
=
\frac{j!}{(r+1)(r+2)\cdots(r+j)}.
\qquad\text{(4)}
\] These identities follow directly from the constant density of the
spacing vector on the simplex.

If \(r\ge d/4\), then \[
\mathbb E(dB_r)^j\le 4^j j!.
\] Consequently, for \(|\theta|\le1/8\), expansion of the exponential
gives \begin{equation*}
\begin{aligned}
\log\mathbb E e^{\theta dB_r}
&\le
\mathbb E e^{\theta dB_r}-1\\
&\le
\theta\frac{d}{r+1}
+\sum_{j\ge2}(4|\theta|)^j\\
&\le
\theta\frac{d}{r+1}+32\theta^2.
\end{aligned}
\qquad\text{(5)}
\end{equation*} The same bound holds for a variable identically zero,
with its mean replacing \(d/(r+1)\).

\begin{center}\rule{0.5\linewidth}{0.5pt}\end{center}

\subsubsection{3. Exposing the half-interval
bits}\label{exposing-the-half-interval-bits}

Generate the labels as \[
X_v=\frac{C_v+U_v}{2},
\] where all \(C_v\sim\operatorname{Bernoulli}(1/2)\) and
\(U_v\sim\mathrm{Unif}[0,1]\) are mutually independent.

Put \[
V_a=\{v:C_v=a\},\qquad a\in\{0,1\}.
\] For each vertex, define its degree across this random cut by \[
r_v=\bigl|\{u\in N_G(v):C_u\ne C_v\}\bigr|,
\] and define \[
b_v=
\begin{cases}
0,&C_v=0,\\
d-r_v,&C_v=1.
\end{cases}
\]

Almost surely, the degree in \(H\) is \[
\deg_H(v)
=
b_v+
\sum_{\substack{u\in N_G(v)\\ C_u\ne C_v}}
\mathbf 1_{\{U_u+U_v\ge1\}}.
\qquad\text{(6)}
\] Indeed, monochromatic edges are deterministic after exposing the
bits, while a cross-edge is retained exactly when \(U_u+U_v\ge1\).

For a fixed bit configuration, let \[
q_{v,k}(C)
=
\frac{\mathbf 1_{\{b_v\le k\le b_v+r_v\}}}{r_v+1},
\qquad
Q_k(C)=\sum_v q_{v,k}(C).
\qquad\text{(7)}
\] The spacing identity, or integration of a binomial distribution,
shows that \[
q_{v,k}(C)=\Pr(\deg_H(v)=k\mid C).
\] Thus \[
Q_k(C)=\mathbb E[m_k\mid C].
\qquad\text{(8)}
\]

We will use the event \[
\mathcal E=\{r_v\ge d/4\text{ for every }v\}.
\] Conditional on either value of \(C_v\), the variable \(r_v\) is
\(\operatorname{Bin}(d,1/2)\). The elementary binomial lower-tail bound
therefore gives \[
\Pr(r_v<d/4)\le e^{-d/16},
\qquad
\Pr(\mathcal E^c)\le ne^{-d/16}.
\qquad\text{(9)}
\]

\begin{center}\rule{0.5\linewidth}{0.5pt}\end{center}

\subsubsection{4. Concentration conditional on the
bits}\label{concentration-conditional-on-the-bits}

Fix a bit configuration belonging to \(\mathcal E\), and fix \(k\).
Define \[
M_{a,k}=\sum_{v\in V_a}\mathbf 1_{\{\deg_H(v)=k\}},
\qquad
Q_{a,k}=\sum_{v\in V_a}q_{v,k}.
\] Thus \(m_k=M_{0,k}+M_{1,k}\) and \(Q_k=Q_{0,k}+Q_{1,k}\).

We first bound the moment-generating function of \(M_{a,k}\).

Condition additionally on all \(U_u\) with \(u\in V_{1-a}\). The degree
indicators for vertices in \(V_a\) are then independent: by (6), each is
a function only of its own remaining label \(U_v\).

Let their conditional success probabilities be \(p_{v,k}\). If \[
0\le k-b_v\le r_v,
\] then \(p_{v,k}\) is the corresponding spacing among the \(r_v\)
points \[
(1-U_u)_{\substack{u\in N_G(v)\\C_u\ne C_v}}.
\] Otherwise \(p_{v,k}=0\). Consequently, \[
\mathbb E[p_{v,k}\mid C]=q_{v,k},
\qquad\text{(10)}
\] and, in the nonzero case, its marginal conditional distribution is
\(\operatorname{Beta}(1,r_v)\).

For real \(\lambda\), set \(\eta=e^\lambda-1\). Conditional independence
gives \begin{equation*}
\begin{aligned}
\mathbb E[e^{\lambda M_{a,k}}\mid C]
&=
\mathbb E_{U_{V_{1-a}}}
\prod_{v\in V_a}(1+\eta p_{v,k})\\
&\le
\mathbb E_{U_{V_{1-a}}}
\exp\left(\eta\sum_{v\in V_a}p_{v,k}\right).
\end{aligned}
\qquad\text{(11)}
\end{equation*} This inequality is valid for either sign of \(\lambda\),
since \(\eta>-1\).

Each opposite-side label is used in at most \(d\) of the functions
\(p_{v,k}\). Applying (2) with \(R=d\), and then (5), yields, for
\(|\eta|\le1/8\), \begin{equation*}
\begin{aligned}
\log\mathbb E[e^{\lambda M_{a,k}}\mid C]
&\le
\frac1d\sum_{v\in V_a}
\log\mathbb E[e^{\eta d p_{v,k}}\mid C]\\
&\le
\eta Q_{a,k}+32\frac{|V_a|}{d}\eta^2.
\end{aligned}
\qquad\text{(12)}
\end{equation*} On \(\mathcal E\), \[
Q_{a,k}\le \frac{4|V_a|}{d}.
\] Since \(e^\lambda-1-\lambda=O(\lambda^2)\), equation (12) implies
that, for absolute constants \(C_0,\lambda_0>0\), \[
\log\mathbb E\!\left[
e^{\lambda(M_{a,k}-Q_{a,k})}\mid C
\right]
\le
C_0\frac{|V_a|}{d}\lambda^2,
\qquad |\lambda|\le\lambda_0.
\qquad\text{(13)}
\]

The two sides need not be independent. Cauchy--Schwarz nevertheless
gives \begin{equation*}
\begin{aligned}
\mathbb E[e^{\lambda(m_k-Q_k)}\mid C]
&\le
\left(
\mathbb E[e^{2\lambda(M_{0,k}-Q_{0,k})}\mid C]
\right)^{1/2}\\
&\quad{}\times
\left(
\mathbb E[e^{2\lambda(M_{1,k}-Q_{1,k})}\mid C]
\right)^{1/2}.
\end{aligned}
\end{equation*} Hence, after adjusting absolute constants, \[
\log\mathbb E[e^{\lambda(m_k-Q_k)}\mid C]
\le C_1\frac nd\lambda^2,
\qquad |\lambda|\le\lambda_1.
\qquad\text{(14)}
\]

Chernoff optimization proves the following estimate, uniformly over
every bit configuration in \(\mathcal E\): \[
\Pr\bigl(|m_k-Q_k|\ge t\mid C\bigr)
\le
2\exp\left[
-c\min\left\{\frac{dt^2}{n},t\right\}
\right].
\qquad\text{(15)}
\]

This handles both tails and all degree values, including \(k=0,d\).

\begin{center}\rule{0.5\linewidth}{0.5pt}\end{center}

\subsubsection{5. Concentration of the random conditional
mean}\label{concentration-of-the-random-conditional-mean}

We next control \(Q_k(C)\). It is important here \textbf{not} to
condition the independent bits on \(\mathcal E\).

Instead, define local truncations \[
\bar q_{v,k}
=
q_{v,k}\mathbf 1_{\{r_v\ge d/4\}},
\qquad
\bar Q_k=\sum_v\bar q_{v,k}.
\] For every bit configuration, \[
0\le\bar q_{v,k}\le\frac4d.
\qquad\text{(16)}
\] Moreover, \(\bar q_{v,k}\) depends only on the bits in \(N_G[v]\).
Each bit belongs to exactly \(d+1\) such closed neighborhoods.

Apply (3) with \(R=d+1\), followed by the elementary Hoeffding
moment-generating-function bound for a variable in an interval of length
\(4/d\). For every real \(\lambda\), \begin{equation*}
\begin{aligned}
\log\mathbb E e^{\lambda(\bar Q_k-\mathbb E\bar Q_k)}
&\le
\frac1{d+1}\sum_v
\log\mathbb E
e^{\lambda(d+1)(\bar q_{v,k}-\mathbb E\bar q_{v,k})}\\
&\le
\frac{n(d+1)}8\left(\frac4d\right)^2\lambda^2\\
&\le
4\frac nd\lambda^2.
\end{aligned}
\end{equation*} Therefore \[
\Pr\bigl(|\bar Q_k-\mathbb E\bar Q_k|\ge t\bigr)
\le
2\exp\left(-\frac{dt^2}{16n}\right).
\qquad\text{(17)}
\]

We also need its mean. Conditional on \(X_v=x\), the original degree is
\(\operatorname{Bin}(d,x)\), so \[
\Pr(\deg_H(v)=k)
=
\int_0^1\binom dk x^k(1-x)^{d-k}\,dx
=
\frac1{d+1}.
\qquad\text{(18)}
\] Thus \(\mathbb Em_k=\mathbb EQ_k=\mu\). Since \(0\le q_{v,k}\le1\),
equations (9) and (18) give \[
0\le
\mu-\mathbb E\bar Q_k
=
\sum_v\mathbb E\!\left[
q_{v,k}\mathbf 1_{\{r_v<d/4\}}
\right]
\le ne^{-d/16}.
\qquad\text{(19)}
\] Finally, on \(\mathcal E\), \[
Q_k=\bar Q_k
\qquad\text{for every }k.
\qquad\text{(20)}
\]

\begin{center}\rule{0.5\linewidth}{0.5pt}\end{center}

\subsubsection{\texorpdfstring{6. Simultaneous concentration for
\(d\ge64\log n\)}{6. Simultaneous concentration for d\textbackslash ge64\textbackslash log n}}\label{simultaneous-concentration-for-dge64log-n}

Put \(B=ne^{-d/16}\). Combining (15), (17), (19), and (20), and taking a
union bound over the \(d+1\) degree values, gives \[
\boxed{
\Pr\left(
\max_{0\le k\le d}|m_k-\mu|\ge 2t+B
\right)
\le
B+
4(d+1)\exp\left[
-c\min\left\{\frac{dt^2}{n},t\right\}
\right].
}
\qquad\text{(21)}
\] Here the bad event \(\mathcal E^c\), which is common to all \(k\), is
counted only once. This inequality is valid for every \(d\ge1\),
although it is useful in this form when \(d\) is at least a sufficiently
large multiple of \(\log n\).

Suppose now that \(d\ge64\log n\). Then \[
B\le n^{-3}.
\] Choose \[
t=A\left(\sqrt{\frac nd\log n}+\log n\right),
\] where \(A\ge1\) is a sufficiently large absolute constant. We have \[
\frac{dt^2}{n}\ge A^2\log n,
\qquad
t\ge A\log n.
\] Thus (21) proves (1).

Under the conjectured hypothesis \(d=o(n/\log n)\), \[
\frac{
\sqrt{(n/d)\log n}+\log n
}{
n/(d+1)
}
=
O\left(
\sqrt{\frac{d\log n}{n}}
+
\frac{d\log n}{n}
\right)
=o(1).
\qquad\text{(22)}
\] This proves Property \((*)\) in the large-degree case.

\begin{center}\rule{0.5\linewidth}{0.5pt}\end{center}

\subsubsection{7. The remaining small
degrees}\label{the-remaining-small-degrees}

For completeness, we handle all \(d<64\log n\) directly.

For fixed \(k\), let \[
I_{v,k}=\mathbf 1_{\{\deg_H(v)=k\}}.
\] Each indicator depends on the independent labels in \(N_G[v]\), so
these functions have the read-\(R\) property with \[
R=d+1.
\] Each indicator is Bernoulli with mean \(p=1/(d+1)\), by (18). Hence
(2) gives \[
\mathbb E e^{\lambda m_k}
\le
\left(1-p+p e^{\lambda R}\right)^{n/R}.
\] The usual elementary Chernoff calculation therefore yields, for
\(0<\varepsilon\le1\), \[
\Pr(|m_k-\mu|\ge\varepsilon\mu)
\le
2\exp\left(
-\frac{\varepsilon^2 n}{3(d+1)^2}
\right).
\qquad\text{(23)}
\] Taking a union bound, \[
\Pr\left(
\max_k|m_k-\mu|\ge\varepsilon\mu
\right)
\le
2(d+1)\exp\left(
-\frac{\varepsilon^2 n}{3(d+1)^2}
\right).
\qquad\text{(24)}
\]

For \(d<64\log n\), choose \[
\varepsilon_n=4(d+1)\sqrt{\frac{\log n}{n}}.
\] Uniformly throughout this degree range, \(\varepsilon_n=o(1)\), and
(24) tends to zero. The case \(d=0\) is deterministic: \(m_0=n\).

Together with Section 6, this covers every sequence \(d=o(n/\log n)\),
completing the proof. \(\square\)

\begin{center}\rule{0.5\linewidth}{0.5pt}\end{center}

\subsubsection{8. Scope and relation to the earlier
obstruction}\label{scope-and-relation-to-the-earlier-obstruction}

The supplied attempt's read-\(R\) argument is valid, but a proper
coloring of \(G\) is unnecessary. Exposing the half-interval bit of
every label makes the \textbf{random part} of the model bipartite,
regardless of \(\chi_f(G)\). The monochromatic edges contribute only the
deterministic offsets \(b_v\).

The remaining conditional mean is controlled because, outside
exponentially unlikely local cut imbalances, each of its summands is at
most \(4/d\) and reads only \(d+1\) independent bits. This removes the
chromatic-number loss and avoids the unresolved residual estimate in the
earlier approach.

There is no remaining mathematical gap for the stated simple-graph
model. The argument does not establish failure beyond the conjectured
range, and literature priority remains unchecked.

\end{document}
